\documentclass[11pt,letterpaper,roman]{moderncv} % Font sizes: 10, 11, or 12; paper sizes: a4paper, letterpaper, a5paper, legalpaper, executivepaper or landscape; font families: sans or roman
%\documentclass[11pt,a4paper,roman]{moderncv}
\moderncvstyle{classic} % CV theme - options include: 'casual' (default), 'classic', 'oldstyle' and 'banking'
\moderncvcolor{blue} % CV color - options include: 'blue' (default), 'orange', 'green', 'red', 'purple', 'grey' and 'black'
\usepackage{textgreek}
\usepackage[
top    = 0.75in,
bottom = 0.75in,
left   = 0.75in,
right  = 0.75in]{geometry}

\usepackage{lastpage}
\rfoot{\addressfont\itshape\textcolor{gray}{\thepage\ / \pageref{LastPage}}}

%\usepackage[scale=0.9]{geometry} % Reduce document margins
\setlength{\hintscolumnwidth}{2.4cm} % Uncomment to change the width of the dates column
%\setlength{\makecvtitlenamewidth}{10cm} % For the 'classic' style, uncomment to adjust the width of the space allocated to your name

%----------------------------------------------------------------------------------------
%	NAME AND CONTACT INFORMATION SECTION
%----------------------------------------------------------------------------------------
%\firstname{\Huge Bingjie} % Your first name 
%\familyname{\Huge Wang} % Your last name
\firstname{Bingjie}
\familyname{Wang}

% All information in this block is optional, comment out any lines you don't need
%\title{Curriculum Vitae {\small (updated 11/2014)}}
%\title{Curriculum Vitae}
%\title{R\ittexttt{\'{e}}sum\ittexttt{\'{e}}}
\email{bwang@jhu.edu}
%\address{3900 N Charles Street}{Baltimore, MD 21218, USA}
\address{Physics and Astronomy, Johns Hopkins University}{3400 N. Charles Street, Baltimore, MD 21218, USA}
%\phone{(000) 111 1112}
%\fax{(000) 111 1113}
%\mobile{+1 (412) 304-6768} 
\renewcommand*\httplink[2][]{{\urlstyle{sf}\expandafter\href#2}}
%\homepage{{http://pages.jh.edu/~bwang50}{pages.jh.edu/\string~bwang50}} %\homepage{staff.org.edu/~jsmith}{staff.org.edu/$\sim$jsmith} % The first argument is the url for the clickable link, the second argument is the url displayed in the template - this allows special characters to be displayed such as the tilde in this example
\homepage{{https://wangbingjie.github.io}{https://wangbingjie.github.io}}
%\extrainfo{Citizenship: P.R. China}
%\photo[55pt][0.0pt]{/Users/sweven/Documents/Applications/pitt} % The first bracket is the picture height, the second is the thickness of the frame around the picture (0pt for no frame)
%\quote{"A witty and playful quotation" - John Smith}

%----------------------------------------------------------------------------------------

\begin{document}

\makecvtitle % Print the CV title

\definecolor{ablue}{rgb}{0.36, 0.54, 0.66}
\definecolor{dgray}{rgb}{0.4, 0.4, 0.4}
%----------------------------------------------------------------------------------------
%	EDUCATION SECTION
%----------------------------------------------------------------------------------------
\vspace{-0.3in}
\section{Research interests}
\cvitem{}{\textbf{Cosmology and extragalactic astronomy}}
\cvitem{}{Interstellar and circumgalactic media, outflow, star-forming galaxy, galaxy evolution, reionization, 21-cm cosmology, the cosmic microwave background}

\section{Education}
%\cvitem{Current}{\textbf{Doctor of Philosophy}, Astrophysics, Johns Hopkins University}
\cvitem{2016 -- 2021}{\textbf{Johns Hopkins University}, Baltimore, MD
\newline Ph.D., Astrophysics, advisor: Prof. Timothy M. Heckman
\newline Thesis: Implications for the epoch of reionization in the local universe
}
\cvitem{2012 -- 2016}{\textbf{University of Pittsburgh}, Pittsburgh, PA
\newline {\emph{Magna Cum Laude}}
\newline B.A., Philosophy
\newline B.Phil., Physics with honors, advisor: Prof. Arthur B. Kosowsky
\newline Thesis: Evaluating the standard model of cosmology in light of large-scale anomalies in the cosmic microwave background
}
%\cvitem{2016}{\textbf{Bachelor of Philosophy in Physics with honors} {\hfill \emph{Magna Cum Laude}}
%\newline University of Pittsburgh \hfill Advisor: Prof. Arthur Kosowsky
%\newline Thesis: Evaluating the standard model of cosmology in light of large-scale anomalies in the cosmic microwave background}
%\cvitem{2016}{\textbf{Bachelor of Arts in Philosophy } \hfill  {\emph{Magna Cum Laude}}
%\newline \hspace{3.6in} University of Pittsburgh}

\section{Professional Experience}
\cvitem{Exp. 2022}{\textbf{The Pennsylvania State University}, University Park, PA
\newline Postdoctoral scholar, mentor: Prof. Joel Leja
}


\section{Refereed Publications}
\vspace{-9pt}
\cvline{}{\color{dgray} Total citations: 144; h-index: 7 {(ADS, 11/01/21)\newline \href{https://ui.adsabs.harvard.edu/search/p_=0&q=orcid\%3A0000-0001-9269-5046&sort=date\%20desc\%2C\%20bibcode\%20desc}{Click here to see all papers in ADS; each title below contains an individual link.}}}
\cvline{Primary- [1] author\hspace{26pt}}{\href{https://ui.adsabs.harvard.edu/abs/2021ApJ...916....3W/abstract}{The Low-redshift Lyman-continuum Survey: [S~\textsc{ii}] deficiency and the leakage of ionizing radiation,}
\newline \textbf{B. Wang}, T. Heckman, and the survey collaboration, ApJ 916, 3 (2021)
\newline -- \href{https://twitter.com/AAS_Office/status/1440419779689410560}{featured in {\textbf{AAS Journal Author Series}}}}
\cvline{[2]}{\href{https://ui.adsabs.harvard.edu/abs/2020ApJ...894..149W/abstract}{A systematic study of galactic outflows via fluorescence emission: implications for their size and structure,}
\newline \textbf{B. Wang}, T. Heckman, G. Zhu, and C. Norman, ApJ 894, 149 (2020)
\newline -- \href{https://aasnova.org/2020/08/07/tracing-gas-flows-out-of-star-forming-galaxies/}{featured in {\textbf{AAS NOVA Research Highlights}}}}
\cvline{[3]}{\href{https://ui.adsabs.harvard.edu/abs/2019ApJ...885...57W/abstract}{A new technique for finding galaxies leaking Lyman-continuum radiation: [S~\textsc{ii}] deficiency,}
\newline \textbf{B. Wang}, T. Heckman, C. Leitherer, et al., ApJ 885, 57 (2019)}
\cvline{[4]}{\href{https://ui.adsabs.harvard.edu/abs/2018ApJ...863..121W/abstract}{A projected estimate of the reionization optical depth using the CLASS experiment's sample variance limited E-mode measurement,}
\newline D. Watts, \textbf{B. Wang}, and the CLASS collaboration, ApJ 863, 121 (2018)
%\newline -- contribution: developed part of the analysis for maximizing the likelihood function
}
\cvline{[5]}
{\href{https://ui.adsabs.harvard.edu/abs/2015PhRvD..92f3008A/abstract}{Microwave background correlations from dipole anisotropy modulation,} 
\newline S. Aiola, \textbf{B. Wang}, A. Kosowsky, et al., PRD 92 (6), 063008 (2015)
%\newline -- contribution: calculated the dipole anisotropy on CMB maps taken by the Planck satellite
}
\cvline{[6]}{\href{https://ui.adsabs.harvard.edu/abs/2015PhRvD..91d3510A/abstract}{Gaussian approximation of peak values in the integrated Sachs-Wolfe effect,} 
\newline S. Aiola, A. Kosowsky, and \textbf{B. Wang}, PRD 91 (4), 043510 (2015)
%\newline -- contribution: calculated peak statistics on simulated CMB maps with desired correlations
}
\cvline{Others \hspace{8pt}  [7]}{\href{https://ui.adsabs.harvard.edu/abs/2022arXiv220111716F/abstract}{The Low-redshift Lyman-continuum Survey I: new, diverse local Lyman-continuum emitters,} \newline S. Flury, et al. (including B. Wang), arXiv:2201.11716 (2022)
}
\cvline{[8]}{\href{https://ui.adsabs.harvard.edu/abs/2019ApJ...876..126A/abstract}{The Baltimore Oriole's Nest: cool winds from the inner and outer parts of a star-forming galaxy at $z=1.3$,} \newline W. Wang, et al. (including B. Wang), arXiv:2109.12133 (2021)
}
\cvline{[9]}{\href{https://ui.adsabs.harvard.edu/abs/2019ApJ...876..126A/abstract}{On-sky performance of the CLASS Q-band telescope,} \newline J. Appel, et al. (including B. Wang), ApJ 876, 126 (2019)
}
\cvline{[10]}{\href{https://ui.adsabs.harvard.edu/abs/2018arXiv181008482K/abstract}{Fermi-LAT counterparts of IceCube neutrinos above 100 TeV,} \newline F. Krau\ss,  et al. (including B. Wang),  A\&A 620, A174  (2018)
%\newline -- contribution: devised broadband spectral energy distributions from Swift and Fermi data
}

%\cvline{Proceedings}{\href{http://arxiv.org/abs/1502.02147}{TANAMI counterparts to IceCube high-energy neutrino events,} \newline
%F. Krau\ss, B. Wang, C. Baxter, M. Kadler, \emph{et al.}, Proceedings of 5\textsuperscript{th} Fermi Symposium (2015), arXiv:1502.02147.}
%\cvline{}{\href{http://dx.doi.org/10.1017/S1743921314013763}{Extreme-value statistics for testing dark energy,} \newline S. Aiola, A. Kosowsky, and B. Wang, Proceedings of the International Astronomical Union, 10 (S306), 54-56, (2014).}
%\newpage

%\section{Publications}
% Publications from a BibTeX file without multibib\renewcommand*{\bibliographyitemlabel}{\@biblabel{\arabic{enumiv}}}% for BibTeX numerical labels
%\nocite{*}
%\bibliographystyle{plain}s
%\bibliography{publications}       % 'publications' is the name of a BibTeX file

%\section{Research projects}
%\cvitem{\color{ablue}$\circ$}{Constraints on cosmological models:\newline
%$\color{gray}\circ$ Investigated anomalies in the standard model of cosmology;\newline
%$\color{gray}\circ$ Explored the parameter space of exotic physics;\newline
%$\color{gray}\circ$ Performed MCMC to maximize likelihood on a computing cluster.}
%
%\cvitem{\color{ablue}$\circ$}{Ionization of the circumgalactic medium:\newline
%$\color{gray}\circ$ Developed a new technique for finding a special class of galaxies;\newline
%$\color{gray}\circ$ Extracted signals from low signal-to-noise observations;\newline
%$\color{gray}\circ$ Searched for correlations in data to infer physical properties of galaxies.}
%
%\cvitem{\color{ablue}$\circ$}{Structure of galactic outflows:\newline
%$\color{gray}\circ$ Devised a novel method to estimate the density and velocity profiles;  \newline
%$\color{gray}\circ$ Analyzed spectroscopic and imaging data;\newline
%$\color{gray}\circ$ Incorporated toy models in radiative transfer simulations.}

% HONORS AWARDS
\section{Honors and Awards}
\cvline{2020}{Rodger Doxsey Prize, American Astronomical Society}
\cvline{2019}{First-prize poster, First Light at University of S\~{a}o Paulo}
\cvline{2016}{$\Sigma \Pi \Sigma$ physics honors society initiate}
\cvline{2012 -- 2016}{Dean's honor list}
\cvline{2016}{Julia Thompson award for excellence in scientific writing, University of Pittsburgh}
\cvline{2015}{Halliday award for excellence in undergraduate research, University of Pittsburgh}
\cvline{2015}{Thomas-Lain fund scholarship, University of Pittsburgh}
\cvline{2014}{Research Internship in Science $\&$ Engineering, Deutschen Akademischen Austauschdienstes}
%\cvline{2014}{DAAD Research Internship in Science $\&$ Engineering \newline
%Administered by the German Academic Exchange Service (DAAD); supported a 12-week summer research internship at a German institute.}
\cvline{2014}{Emil Sanielevici undergraduate research scholarship, University of Pittsburgh}
%\cvline{2013}{The National Society of Collegiate Scholars Invitee}
%\cvline{2013}{The National Society of Leadership and Success, $\Sigma$A$\Pi$, Invitee}
%\cvline{2013}{$\Delta$EI Academic Honor Society Invitee}

\section{Technical Skills}
\cvline{Proficient}{
%\vspace{-9pt}{\color{ablue}$\circ$} 
Python, C, Git, \LaTeX;
\newline 
%{\color{ablue}$\circ$}  
Bayesian analysis, the Cosmic Linear Anisotropy Solving System (CLASS) \& modifications, HEALPix, HST UV spectroscopic \& imaging analyses, high-performance computing, MCMC, PolSpice.
}
\cvline{Working knowledge}{
%\vspace{-9pt}{\color{ablue}$\circ$} 
Java, SQL, Fortran90, Bash/shell, Mathematica, HTML;
\newline
%{\color{ablue}$\circ$} 
21cmFAST, Radiation Scattering in Astrophysical Simulations (RASCAS), machine learning.
}

%\section{Media coverage}
%\cvline{08/07/20}{\href{https://aasnova.org/2020/08/07/tracing-gas-flows-out-of-star-forming-galaxies/}{Tracing gas flows out of star-forming galaxies}
%\newline
%by Tarini Konchady in AAS NOVA research highlights.}

\section{Presentations}
%\cvitem{08/27/21}{Invited talk, AAS Journal Author Series, \newline
%{} "The Low-redshift Lyman-continuum Survey: [S~\textsc{ii}] deficiency and the leakage of ionizing radiation," \newline
%To be publicly available mid September.
%}
\cvitem{03/12/21}{Invited talk, STARs at Arizona State University, \newline
{} "Star-forming galaxies: ionizing-photon escape and outflow scale"}
\cvitem{01/12/21}{Dissertation talk, 237$^{\textrm{th}}$ Meeting of the American Astronomical Society, \newline
{} "Implications for the epoch of reionization in the local universe"}
\cvitem{10/29/20}{Department lunch talk, University of California at Berkeley, \newline
{} "Implications for the epoch of reionization in the local universe"}
\cvitem{02/24/20}{CAS wine \& cheese, Johns Hopkins University, \newline
{} "A systematic study of galactic outflows via fluorescence emission: implications for their size and structure"}
\cvitem{08/06/19}{First light, University of S\~{a}o Paulo, Brazil, \newline
{} "A new technique for finding galaxies leaking Lyman-continuum radiation: [S~\textsc{ii}] deficiency"}
\cvitem{10/22/18}{CAS wine \& cheese, Johns Hopkins University, \newline
{} "Modeling the global 21-cm signal"}
\cvitem{08/06/15}{Undergraduate summer research, Princeton University, \newline
{} "Filtering the microwave background temperature through correlation with polarization"}
\cvitem{02/25/15}{Annual Sanielevici lecture, University of Pittsburgh, \newline
{} "Re-examining dark energy through large-scale structures of the universe"}
\cvitem{09/12/14}{Workshop on large-scale anomalies, Case Western Reserve University, \newline
{} "Gaussian approximation of peak values in the integrated Sachs-Wolfe effect."}
\cvitem{07/05/14}{DAAD RISE scholarship holder meeting, Heidelberg, Germany, \newline
{} "Black holes as possible sources of high-energy neutrinos"}
\cvitem{04/03/14}{Workshop on astrophysics and cosmology, Pennsylvania State University, \newline
{} "$\Lambda$CDM estimates on the stacked late-ISW signal"}

%\newpage
%\section{Research Experience}
%\cvitem{Current}{Graduate research assistant, Johns Hopkins University}
%\cvitem{}{Advisor: Prof. Timothy Heckman}
%\cvitem{}{Ionization of the intergalactic medium and starburst-driven outflows.}
%\cvitem{09/16 - 04/18}{Graduate research assistant, Johns Hopkins University}
%\cvitem{}{Advisor: Prof. Tobias Marriage}
%\cvitem{}{Collaboration: The Cosmology Large Angular Scale Surveyor (CLASS)}
%\cvitem{}{Likelihood analysis to extract cosmological parameters from simulated CLASS maps.}
%\cvitem{05/13 - 04/16}{Undergraduate research assistant, University of Pittsburgh}
%\cvitem{}{Advisor: Prof. Arthur Kosowsky}
%\cvitem{}{Large-scale anomalies in the Cosmic Microwave Background: i. hemispherical power asymmetry; ii. excess signal of the integrated Sachs-Wolfe effect.}
%\cvitem{06/15 - 08/15}{Undergraduate summer research, Princeton University}
%\cvitem{}{Supervisor: Prof. Ren\'ee Hlo{\v z}ek (now faculty at Dunlap Institute, University of Toronto)}
%\cvitem{}{{Two weeks of summer school focusing on data analysis and computational science and one week on high performance computing, followed by directed research on cross-correlation between CMB temperature and polarization.}}
%\cvitem{05/14 - 08/14}{DAAD RISE, Dr. Karl Remeis Observatory, Bamberg, Germany}
%\cvitem{}{Supervisors: Prof. J\"orn Wilms \& Dr. Felicia Krau\ss}
%\cvitem{}{{Active Galactic Nuclei as candidates for the origin of a high-energy neutrino flux detected by IceCube.}}

%\newpage
\section{Community Service}
\cvitem{}{Reviewer for The Astrophysical Journal}
\cvitem{}{Volunteer at Physics Fair, Baltimore, MD}

\section{Teaching Experience}
\cvitem{09/16 - 12/18}{Graduate teaching assistant; courses taught:}
\cvitem{}{
Graduate radiative astrophysics; graduate astrophysical dynamics \hfill Fall 2018 \newline
Graduate radiative astrophysics; cosmology \hfill Fall 2017 \newline
General physics for biological science majors II; general physics lab II \hfill Spring 2017 \newline
General physics for physical science majors II; general physics lab I \hfill Fall 2016}
%\vspace{0.05in}
%\cvitem{01/15 - 04/15}{Undergraduate teaching assistant for physics resource room}
%\vspace{0.05in}
%\cvitem{01/13 - 04/13}{Undergraduate teaching assistant for introduction to physics I}

%\newpage
\section{Summer Schools}
\cvitem{07/08 - 08/07/19}{First light: stars, galaxies and black holes in the epoch of reionization, \newline University of S\~{a}o Paulo, Brazil.} 
\cvitem{07/03 - 28/17}{Particles, strings and cosmology, University of Hamburg, Germany.}

%\section{Other Conferences Attended}
%\cvline{05/15}{Phenomenology 2015 Symposium, \newline
%University of Pittsburgh, Pittsburgh, PA, USA.}
%\cvline{02/15}{The Atacama Cosmology Telescope collaboration meeting, \newline
%Princeton University, Princeton, NJ, USA.}
%%\cvline{}{The Atacama Cosmology Telescope collaboration meeting}
%\cvline{05/14}{Phenomenology 2014 Symposium, \newline
%University of Pittsburgh, Pittsburgh, PA, USA.}

\section{References}
\cvline{}{Timothy M. Heckman, Ph.D. \hspace{97pt}  Ely D. Kovetz, Ph.D.\newline 
Chair and A. Hermann Pfund Professor \hspace{47pt} Senior Lecturer\newline
Department of Physics \& Astronomy \hspace{60pt} Department of Physics\newline
Johns Hopkins University \hspace{115pt} Ben-Gurion University\newline
theckma1@jhu.edu \hspace{149pt} kovetz@bgu.ac.il\newline
+1 410-516-7369, 410-516-5454 \hspace{89pt} +972 64-61167}
\vspace{0.05in}
\cvline{}{Tobias A. Marriage, Ph.D. \hspace{110pt} Arthur B. Kosowsky, Ph.D. \newline
Associate Professor \hspace{145pt} Chair and Professor\newline
Department of Physics \& Astronomy \hspace{60pt} Department of Physics \& Astronomy\newline
Johns Hopkins University \hspace{115pt} University of Pittsburgh\newline
marriage@jhu.edu \hspace{153pt} kosowsky@pitt.edu \newline
+1 410-516-6526}

\end{document}

\cvitem{\color{ablue}\textbullet}{Research works focused on large-scale anomalies, in particular: (see publications for details)
\newline
\vspace{-0.17in}
\begin{itemize}
\item Excess signal of Integrated Sachs-Wolfe Effect
\begin{itemize}
\item worked with WMAP $\&$ Planck data sets;
\end{itemize}
\item Hemispherical power asymmetry.
%\end{itemize}
\begin{itemize}
\item worked with WMAP $\&$ Planck data sets;
\item wrote various Python codes for analytical/numerical analyses; etc.
\end{itemize}
\end{itemize}}